\documentclass[xcolor=table,10pt,t]{beamer}
\title{FLAT}
\subtitle{FoLiA Linguistic Annotation Tool}
\date{\today}
\author{Maarten van Gompel}
%% Use the official themes as much as possible
%\usetheme[official=false,department=clst]{ruhuisstijl}
%\usetheme[official=false]{ruhuisstijl}
\usetheme[official=true,department=clst]{ruhuisstijl}


\begin{document}

\begin{frame}
  \titlepage
        \begin{figure}
          \includegraphics[height=2cm]{logo.png}
        \end{figure}
        \vspace{2.5cm}
      \includegraphics[justify=left,width=5cm]{lama-logo-transparent.png}
\end{frame}







\begin{frame}{Introduction (1/2)}
  \begin{block}{What is FLAT?}
      \begin{itemize}
        \item \textbf{F}oLiA \textbf{L}inguistic \textbf{A}nnotation \textbf{T}ool
        \item A web-based tool for linguistic annotation of language resources
        \item Many different annotation types
        \item Flat is anything but ``flat'', it deals with deeply nested data
          and retains document structure.
        \item Aims to be a fully-featured annotation environment for FoLiA
          documents (currently in an early beta stage!)
      \end{itemize}
  \end{block}
\end{frame}

\begin{frame}{Introduction (2/2)}
  \begin{block}{What is FoLiA?}
    \begin{itemize}
      \item \textbf{F}ormat for \textbf{Li}nguistic \textbf{A}nnotation
      \item An XML-based file format for storing and exchanging language
        resources
      \item Generic paradigm independent of language or label set
      \item Single format for all
      \item Implemented in many tools \emph{(ucto,frog,
        ucto,pynlpl,ticcl,folia-tools)}
      \item Used in many projects \emph{(SoNaR500, Nederlab, DutchSemCor, Basilex,
        OpenSoNaR)}
      \item Funded by CLARIN-NL
    \end{itemize}
  \end{block}
\end{frame}

\begin{frame}{Notable Features}
  \begin{block}{ }
    \begin{itemize}
      \item Modern \textbf{web-based} interface
      \item \textbf{Multi-user} environment
      \item Download and upload FoLiA XML documents
      \item \textbf{Concurrency} (multiple users may edit the same document similtaneously)
      \item Storage model with [git] \textbf{versioning control} (unlimited undo operations)
      \item Full \textbf{annotator information}: who annotated what and when?
      \item Layered, retains \textbf{document structure}
      \item Annotations may pertain to single tokens or span over multiple
        (span annotation)
      \item Support for \textbf{corrections} of either text or linguistic annotations. Corrections are never mere substitutes, originals are always preserved!
      \item Support for alternative annotations.
      \item Introduces FoLiA Query Language for automated and batch editing
    \end{itemize}
  \end{block}
\end{frame}

\begin{frame}{Visualisation of annotations on hover}
  \begin{figure}
        \includegraphics[height=6cm]{hover.png}
    \end{figure}
\end{frame}

\begin{frame}{Global visualisation of classes}
  \begin{figure}
        \includegraphics[width=15cm]{highlight2.png}
    \end{figure}
\end{frame}

\begin{frame}{Editor}
  \begin{figure}
        \includegraphics[height=6cm]{edit2.png}
    \end{figure}
\end{frame}

\begin{frame}{In the field}
  \begin{block}{Currently used in two projects:}
    \begin{enumerate}
      \item A project with \emph{College van Examens}: Spelling correction on
        texts of children and on automatic corrections by Valkuil.net,
        annotated by four student assistants. 
      \item A project concerning the annotation of named entities, specifically
        locations, in historical texts from World War II, with \emph{Susan Evans},
        \emph{Marten Düring} and \emph{Antal van den Bosch}
    \end{enumerate}
  \end{block}
\end{frame}

\begin{tussenpagina}{ }{ }{ner.png}
\end{tussenpagina}

\begin{frame}{Demo}
  \begin{block}{}
      \begin{itemize}
        \item Demo: \url{http://flat.science.ru.nl}
        \item Source: \url{https://github.com/proycon/flat}
        \item FoLiA website: \url{http://proycon.github.io/folia}
      \end{itemize}

      \medskip
        \begin{figure}
          \includegraphics[height=3cm]{logo.png}
        \end{figure}
  \end{block}
\end{frame}



\end{document}
